\section{3. Variational limit and the critical profile}\label{critical-cost-and-recovery-under-full-marginal-constraints}

Under the traded-coordinate condition of Proposition~\ref{prop:nominal-benchmark}, a local exchange attains the initial slope. We now hold \(s=\rho/\varepsilon^2\) fixed to determine how far that efficiency persists toward reversal.
The limiting objects introduced below connect the stress profile to its reversal coefficient:
\[
\varepsilon^{-2}D_\varepsilon(s\varepsilon^2)
\xrightarrow{\text{locally uniformly}}
F(s)
\xrightarrow{\text{unique zero}}
s^*.
\]
The first passage follows from Theorem~\ref{thm:critical-limit}, which combines convergence of the transport energies with exact marginal recovery and convergence of the paired risk. Strict growth in Theorem~\ref{thm:strict-crossing} then identifies the common critical coefficient.

\subsection{3.1 Local transport and marginal compensation}\label{sec:limit-geometry}
An \(O(\varepsilon)\)-layer around \(\{B=v\}\) contains \(O(\varepsilon)\) reference mass, so displacements of order \(\varepsilon\) change both the transport cost and the paired hinge expectation at order \(\varepsilon^2\). We resolve this layer through the normal coordinate \(u=(B(X)-v)/\varepsilon\) and displacement \(d=(Y-X)/\varepsilon\), while retaining the cost of any transport outside bounded windows needed to restore the marginals.

\paragraph{Definition (Marginal flux and effective energy)} On bounded \(u\)-windows the rescaled source converges to \(d\nu=d\sigma\,du\). We disintegrate a limiting plan with this source into a probability kernel \(\alpha_{\xi,u}(dd)\), specifying the rescaled displacement at surface point \(\xi\) and normal coordinate \(u\); stationary sources have kernel \(\delta_0\). Writing \(d\Pi_\alpha(\xi,u,d)=\alpha_{\xi,u}(dd)\,d\nu(\xi,u)\), we require a finite weighted displacement moment even though the source measure \(\nu=\sigma\otimes du\) has infinite total mass:
\[
I(\alpha)=\int|d|_\gamma\,d\Pi_\alpha<\infty,
\qquad |d|_\gamma=\sum_i\gamma_i|d_i|.
\]
Across a coordinate cut \(\{X_i\le q\}\), upward movements remove mass from below the cut and downward movements add mass to it. The resulting signed flux has a Lebesgue density \(j_i\in L^1(I_i)\), as shown by the section formulas in Appendix A.3, and we define
\begin{equation}
j_i(q)dq=(\xi_i)_\#(d_i\,d\Pi_\alpha),\qquad
\mathcal C(\alpha)=I(\alpha)+\sum_{i=1}^3\gamma_i\|j_i\|_{L^1(I_i)}.
\tag{3.1}\label{eq:legacy-4.1}
\end{equation}
In words, \(j_i\) is the signed displacement moment projected onto the \(i\)-th coordinate, and its negative is the leading marginal CDF discrepancy produced by a regular local plan. Because the full plan preserves each marginal, transport outside a bounded normal and displacement window must compensate the discrepancy within it; the one-dimensional CDF formula for \(W_1\) \cite{vallender1974} gives the resulting cost lower bound. The energy \(\mathcal C\) includes this limiting restoration cost in addition to \(I\) (Figure~\ref{fig:tail-layer}). Projection contracts total variation, giving \(I\le\mathcal C\le2I\), with \(\mathcal C=I\) when all projected fluxes vanish and \(\mathcal C=2I\) when each \(d_i\) has a fixed sign \(\Pi_\alpha\)-almost everywhere, so that \(\|j_i\|_1=\int|d_i|\,d\Pi_\alpha\).

\begin{figure}[tbp]
\centering
\includegraphics[width=\linewidth]{figures/tail-layer-repair.pdf}
\caption{Local transport and full marginal compensation (schematic). A layer with thickness and mass of order \(O(\varepsilon)\), transported over distances of order \(O(\varepsilon)\), incurs second-order cost. The remaining nonnegative plan must cancel the local plan's coordinate CDF discrepancies. Their integrated absolute values give the compensation lower bound, with leading coefficients \(\|j_i\|_1\); a zero flux requires no corresponding leading-order compensation. The shaded discrepancy illustrates this mechanism.}
\label{fig:tail-layer}
\end{figure}


\subsection{3.2 The limiting risk and convergence theorem}\label{sec:limit-risk}
Under the same local plan, each loss has its own minimizing threshold. At threshold \(v+\varepsilon\tau\), a displacement \(d\) changes the leading centered and rescaled loss of \(L_{c,\varepsilon}=B+c\varepsilon H\) from \(u+cH\) to \(u+b\cdot d+cH\). Denote the corresponding hinge increment by \(\mathfrak q_c(\tau)=(u+b\cdot d+cH-\tau)_+-(u+cH-\tau)_+\). We write \(T_c^\alpha(\tau)=\int\mathfrak q_c(\tau)\,d\Pi_\alpha\) for its integral.

\paragraph{Definition (Limiting paired risk)} Moving the threshold away from \(v\) contributes the curvature term \(g\tau^2/2\) from \eqref{eq:legacy-B.2}. Adding the local transport increment gives the two limiting threshold objectives, which must still be minimized separately:
\begin{equation}
G_c^\alpha(\tau)=\frac g2\tau^2+T_c^\alpha(\tau),\qquad
\mathcal V(\alpha)=-\frac Jp+
\frac{\min_{\tau_B}G_0^\alpha(\tau_B)-\min_{\tau_A}G_1^\alpha(\tau_A)}p.
\tag{3.2}\label{eq:legacy-4.2}
\end{equation}

The energy \(\mathcal C\) depends only on the absolute displacement moment and the signed coordinate projections, and is convex under mixtures of kernels. The paired risk also depends on how the displacements are distributed relative to the normal coordinate \(u\): each of the two hinge objectives must optimize its own threshold, so their difference need not be concave in the common kernel.

We omit the superscript \(\alpha\) when the kernel is clear, and write \(\kappa_B=\max_i b_i/\gamma_i\), \(\kappa_H=\max\{\gamma_1^{-1},\gamma_2^{-1}\}\). Since \(|T_c^\alpha|\le\kappa_BI\), the quadratic term gives coercivity and confines every minimizer to \(g\tau^2/2\le2\kappa_BI\). Together with convexity from Lemma~\ref{lem:threshold-convexity}, these properties justify the threshold minimizations in \(\mathcal V\). We then optimize the paired value over the energy budget, with attainment established in Theorem~\ref{thm:critical-limit}:

\begin{equation}
F(s)=\max_{\mathcal C(\alpha)\le s}\mathcal V(\alpha),\qquad s\ge0.
\tag{3.3}\label{eq:legacy-4.3}
\end{equation}

We use vague convergence on \(\mathscr X=\mathcal M^+_{\rm loc}(Z)\), testing against compactly supported continuous functions on \(Z=I^2\times\mathbb R\times\mathbb R^3\), so that compensation may occur outside bounded windows. We lift a plan \(\pi\) from \(P_0\) to a law in \(\Pi\) by
\[
\mathscr T_\varepsilon\pi=\varepsilon^{-1}
\left(X_1,X_2,\frac{B(X)-v}{\varepsilon},\frac{Y-X}{\varepsilon}\right)_\#\pi.
\]
Here the factor \(\varepsilon^{-1}\) compensates for the thickness of the source layer. Because the lifted coordinates determine \((X,Y)\) for fixed \(\varepsilon\), we can retain the cost of the entire original plan in its rescaled energy:

\[
\begin{aligned}
\mathcal E_\varepsilon(\Lambda)
&=\begin{cases}
\displaystyle\int |d|_\gamma\,d\Lambda
=\varepsilon^{-2}\int c_\gamma\,d\pi,
&\begin{gathered}\Lambda=\mathscr T_\varepsilon\pi,\\ \pi\in\Gamma(P_0,P),\ P\in\Pi,\end{gathered}\\
+\infty,&\text{otherwise},
\end{cases}\\
\mathcal E(\Lambda)
&=\begin{cases}
\mathcal C(\alpha),&\Lambda=\Pi_\alpha,\ I(\alpha)<\infty,\\
+\infty,&\text{otherwise}.
\end{cases}
\end{aligned}
\]
Here we identify \(\Pi_\alpha\) with a measure on \(Z\) through \((\xi,u,d)\mapsto(\xi_1,\xi_2,u,d)\), using \(\xi=(x,y,z_v(x,y))\) on \(\Omega\) and assigning zero measure to \((I^2\setminus\Omega)\times\mathbb R\times\mathbb R^3\). Because \(\mathcal E_\varepsilon\) includes the whole plan, including transport outside bounded normal windows, minimizing it over plans with target law \(P\) gives \(\varepsilon^{-2}W_\gamma(P,P_0)\). We formulate the following \(\Gamma\)-convergence statement in this vague topology \cite{dalmaso1993}.

\keepproofstart{7\baselineskip}
\resultnumber{3.1}\begin{theorem}[\(\Gamma\)-convergence and the critical risk profile]\label{thm:critical-limit} Consider the exchangeable box model of \mbox{\hyperref[model:critical-box]{Section~2.2}}.
\begin{resultparts}
\item \emph{Transport plans and compactness.} \(\mathcal E_\varepsilon\) \(\Gamma\)-converges to \(\mathcal E\) as \(\varepsilon\downarrow0\) in the vague topology on \(\mathscr X\). Every sequence \(\varepsilon_n\downarrow0\), \(\Pi_n\in\mathscr X\), with bounded \(\mathcal E_{\varepsilon_n}(\Pi_n)\) has a vaguely convergent subsequence. If \(\Pi_n=\mathscr T_{\varepsilon_n}\pi_n\to\Pi_\alpha\) has bounded energy and \(P_n\) is the target law of \(\pi_n\), then \(\Phi_{\varepsilon_n}(P_n)/\varepsilon_n^2\to\mathcal V(\alpha)\).

\item \emph{Target laws and exact recovery.} For every sequence \(\varepsilon_n\downarrow0\) and every feasible sequence satisfying \(W_\gamma(P_n,P_0)=O(\varepsilon_n^2)\), there are a subsequence and a kernel with finite \(I\) such that

\begin{equation}
\mathcal C(\alpha)\le\liminf_n\frac{W_\gamma(P_n,P_0)}{\varepsilon_n^2},\qquad
\frac{\Phi_{\varepsilon_n}(P_n)}{\varepsilon_n^2}\longrightarrow\mathcal V(\alpha).
\tag{3.4}\label{eq:legacy-4.5}
\end{equation}

Conversely, every kernel with finite \(I\) admits a recovery sequence of distributions preserving the original three full marginals, with

\begin{equation}
\limsup_{\varepsilon\downarrow0}\frac{W_\gamma(P_\varepsilon,P_0)}{\varepsilon^2}\le\mathcal C(\alpha),\qquad
\frac{\Phi_\varepsilon(P_\varepsilon)}{\varepsilon^2}\longrightarrow\mathcal V(\alpha).
\tag{3.5}\label{eq:legacy-4.6}
\end{equation}

\item \emph{The stress profile.} The maximum in \eqref{eq:legacy-4.3} is attained. The function \(F\) is finite, continuous, and nondecreasing, with \(F(0)=-J/p\) and \(F>0\) at some finite budget. For every finite \(S_0\),

\begin{equation}
\sup_{0\le s\le S_0}\left|\frac{D_\varepsilon(s\varepsilon^2)}{\varepsilon^2}-F(s)\right|\longrightarrow0.
\tag{3.6}\label{eq:legacy-4.4}
\end{equation}
\end{resultparts}
\end{theorem}

\paragraph{Proof strategy}
Theorem~\ref{thm:critical-limit}(i) gives \(\Gamma\)-convergence of the transport energy and convergence of the paired risk along bounded-energy plans; exact marginal recovery in part~(ii) makes the limiting kernel values asymptotically attainable in \(\Pi\). Part~(iii) gives uniform convergence of the optimized profiles on compact budget intervals, which together with the strict crossing in Theorem~\ref{thm:strict-crossing} locates the reversal budget \(s^*\).

\emph{Step 1: Localization and the energy lower bound.} Because the two ES thresholds are optimized independently, comparison of the objectives requires a common threshold window for stresses with bounded rescaled energy. In Appendix~\hyperref[proof:threshold-localization]{A.2}, Lemma~\ref{b-l2-all-optimal-thresholds-under-a-budget-of-order-varepsilon2} places every minimizing threshold in an \(O_M(\varepsilon)\) window. On bounded windows the lifted plans are locally compact; outside them, marginal conservation requires compensation for the projected local flux. The compensation bound above then gives the energy lower bound for arbitrary admissible plans. Choosing optimal plans yields the corresponding Wasserstein bound for their target laws.

\emph{Step 2: Convergence of the paired risk.} Bounded first moments make each hinge increment integrable, but need not give uniform integrability along convergent plans. The two objectives share the same large-displacement contribution, so we subtract \(\mathfrak q_0(0)\) from both. The relative increments are bounded on the common threshold window. Outside a normal window enlarged to account for the bounded reallocation \(H\), a nonzero relative increment requires a sufficiently large displacement, so the first moment controls the far-field contributions. The baseline integral is finite but need not converge. Since it translates both threshold objectives by the same amount independently of the threshold, it cancels exactly from their difference of minima.

\emph{Step 3: Exact marginal recovery.} Let \(\mathcal K_{\rm fin}^c\) denote probability kernels with compactly supported continuous weights on finitely many fixed displacement vectors and the remaining mass at zero, as in \eqref{eq:legacy-B.26}. For a fixed kernel in this class, local translations yield a preliminary law \(Q_\varepsilon\) with uniform marginal CDF expansion \(F_{i,\varepsilon}-F_i=-\varepsilon^2j_i+o(\varepsilon^2)\), where \(F_i\) is the reference CDF. The quantile maps \(F_i^{-1}\circ F_{i,\varepsilon}\) restore all three marginals at cost \(\varepsilon^2\sum_i\gamma_i\|j_i\|_1+o(\varepsilon^2)\), with \(O_\alpha(\varepsilon^2)\) displacements.

To see why this correction does not alter the leading paired risk, compare its displacement scale with the \(O_\alpha(\varepsilon)\) threshold band. Outside a sufficiently wide band, the paired hinge has gradient \(-\varepsilon h\) above both thresholds and zero below both, so multiplying by the restoration displacement gives an \(O_\alpha(\varepsilon^3)\) contribution. Inside the band the gradient is bounded, but the \(Q_\varepsilon\)-mass is only \(O_\alpha(\varepsilon)\), giving the same order after integration. The restoration error is therefore \(O_\alpha(\varepsilon^3)\), uniformly over all minimizing threshold pairs.

Approximation by \(\mathcal K_{\rm fin}^c\) in cost, paired risk, and the vague plan topology (Appendix A.4), followed by diagonal recovery, gives the general \(o(\varepsilon^2)\) error. Compactness, contraction toward the identity, and the common budget-Lipschitz bound then give locally uniform convergence of the optimized profiles.

